

\textbf{Reinforcement Learning.}
We formulate the task by defining a Markov Decision Process (MDP) \cite{sutton1998reinforcement} with a tuple $(\mathcal{S}, \mathcal{A}, P, R, \gamma)$.
Here, $\mathcal{S}$ and $\mathcal{A}$ denote the state and action spaces.
We use $o \in \mathcal{S}$ to denote a state and $a \in \mathcal{A}$ to denote an action.
The transition kernel is denoted by $P(o' | o,a)$, the reward function by $R(o,a)$,
and $\gamma \in [0,1)$ is the discount factor.
A policy $\pi(a | o)$ induces a distribution over trajectories
$\tau = (o_0,a_0,o_1,a_1,\ldots)$, where $a_t \sim \pi(\cdot | o_t)$
and $o_{t+1} \sim P(\cdot | o_t,a_t)$.
The standard objective of RL is to maximize
\begin{align}
\label{eq:rl_objective_prelim}
    J(\theta)
    =
    \E_{\tau \sim d_{\pi_\theta}(\tau)}
    \Big[
    \sum_{t \ge 0} \gamma^t R(o_t,a_t)
\Big],
\end{align}
where $d_{\pi_\theta}(\tau)$ is the distribution of trajectories generated by the policy $\pi_\theta$ and the transition kernel.
Policy gradient methods (e.g., REINFORCE \cite{williams1992simple}) update $\theta$ by ascending an estimate of $\nabla_\theta J(\theta)$,
\begin{align}
\label{eq:policy_gradient}
    \nabla_\theta J(\theta)
    =
    \E_{\tau \sim d_{\pi_\theta}(\tau)}
    \Big[\sum_{t\ge 0} \gamma^t
        \nabla_\theta \log \pi_\theta(a_t | o_t)
        \sum_{k\ge t}\gamma^{k-t} R(o_k,a_k)
    \Big],
\end{align}
where $\nabla_\theta \log \pi_\theta(a_t | o_t)$ is the score function of the policy.
To reduce variance, we often replace the cumulative future reward $\sum_{k\ge t}\gamma^{k-t} R(o_k,a_k)$ with an advantage estimate $A(o_t,a_t)$.

In practice, directly optimizing \eqref{eq:policy_gradient}
may be unstable due to overly large policy updates.
Trust-region style policy gradient methods \cite{schulman2015trust,schulman2017proximal,shao2024deepseekmath} instead constrain the policy change between successive iterates.
A popular example is Proximal Policy Optimization (PPO) \cite{schulman2017proximal}, which maximizes $J_{\mathrm{PPO}}(\theta)$
\begin{align}
    \E_{\tau \sim d_{\pi_{\text{old}}}(\tau)}
    \Big[
        \sum_{t\ge 0}
        \min\Big(
            r_\theta(o_t,a_t)\hat{A}_{\mathrm{old}}(o_t,a_t),
            \mathrm{clip}\big(r_\theta(o_t,a_t),1-\epsilon,
            1+\epsilon\big)\hat{A}_{\mathrm{old}}(o_t,a_t)
        \Big)
    \Big],
    \label{eq:ppo}
\end{align}
where $\pi_{\mathrm{old}}$ is the behavior policy, $r_\theta(o,a):=\pi_\theta(a|o)/\pi_{\mathrm{old}}(a|o)$ is the probability ratio, $\epsilon>0$ is the clipping parameter, and $\hat{A}_{\mathrm{old}}(o,a)$ is an advantage estimator.
The clipped surrogate upweights actions that are estimated to be better than baseline and limits the policy update away from $\pi_{\mathrm{old}}$.

\begin{table*}[t]
\centering
\caption{\small Doob Policy Optimization}
\label{tab:unified_framework}
\resizebox{0.8\linewidth}{!}{%
\begin{tabular}{l|ccc}
\toprule
Method
& Generative policy
& Policy update method
& Ratio
\\
\midrule
DPMD \cite{ma2025efficient}
& Diffusion
& Policy Mirror Descent (KL-proximal)
& Analytic
\\

SDAC \cite{ma2025efficient}
& Diffusion
& Soft Actor Critic
& Analytic
\\

DoobPMD (ours)
& Flow
& Policy Mirror Descent (KL-proximal)
& Analytic
\\

DoobPPO (ours)
& Diffusion / Flow
& Proximal Policy Optimization
& Learned
\\

DoobGRPO (ours)
& Diffusion / Flow
& Group Relative Policy Optimization
& Learned
\\
\bottomrule
\end{tabular}}
% \vspace{-2em}
\end{table*}




\textbf{Diffusion Models and Score Matching.}
Let $p_0$ denote the data distribution.
Diffusion models define a forward noising process
$(Y_t)_{t \in [0,T]}$ that gradually transforms data
into noise through the Stochastic Differential Equation
(SDE) \cite{song2020score}
\begin{align}
\label{eq:forward_sde}
    \dd Y_t
    =
    f(Y_t,t)\dd t
    +
    \sigma_t\dd W_t,
    \quad
    Y_0 \sim p_0,
    \quad
    t \in [0,T],
\end{align}
where $f$ and $\sigma_t$ are pre-designed drift and
diffusion coefficients, and $W_t$ is a standard
Brownian motion.
We denote by $p_t$ the marginal distribution of $Y_t$.
The generative process is the reverse-time denoising
process $(X_t)_{t \in [0,T]}$, evolving from $t=T$
to $t=0$.
It transports noise back to data, with
$X_T \sim p_T$ and $X_0 \sim p_0$.
Its time-$t$ marginal is $p_t$, and the process
is governed by the reverse-time SDE
\begin{align}
\label{eq:reverse_sde}
    \dd X_t
    =
    \bigl[
        f(X_t,t)
        -
        \sigma_t^2 \nabla_x \log p_t(X_t)
    \bigr]\dd t
    +
    \sigma_t \dd \overline W_t,
    \quad
    X_T \sim p_T,
    \quad
    t:T \to 0.
\end{align}
Here, $\overline W_t$ denotes a Wiener process in
the reverse-time direction.

Directly using \eqref{eq:reverse_sde} is difficult
in practice because the noise distribution $p_T$
and the true score $\nabla_x \log p_t$ are not
known explicitly.
Therefore, \eqref{eq:reverse_sde} is implemented
through two approximations.
First, the noise distribution $p_T$ is approximated
by a standard Gaussian distribution.
Second, the ground-truth score function
$s(x,t) := \nabla_x \log p_t(x)$
is replaced by a parameterized estimator
$s_\phi(x,t)$.
The score estimator is learned by minimizing
the score matching loss
\begin{align}
\label{eq:score_matching}
    \mathcal{L}_{\mathrm{SM}}(\phi)
    =
    \E_{t, Y_t \sim p_t}
    \bigl[
        \|s(Y_t,t) - s_\phi(Y_t,t)\|_2^2
    \bigr].
\end{align}


\textbf{Flow Matching.}
Following the notation in \cite{lipman2024flow}, let $p_0$ denote the noise distribution and $p_T$ the data distribution (opposite to the diffusion convention).
Flow matching constructs a deterministic transport from $p_0$ to $p_T$ via the Ordinary Differential Equation (ODE)
\begin{align}
\label{eq:fm_ode}
    \dd X_t = u(X_t,t)\dd t,
\end{align}
where $u(x,t)$ is the velocity field and $p_t$ denotes the marginal probability density function of $X_t$.
Given a reference velocity $u(x,t)$ that induces the desired path $\{p_t\}_{t\in[0,T]}$, flow matching learns a parameterized velocity $u_\phi(x,t)$ by minimizing
the flow matching loss
\begin{align}
\label{eq:flow_matching}
    \mathcal{L}_{\mathrm{FM}}(\phi)
    =
    \E_{t, X_t \sim p_t}
    \Big[\big\| u(X_t,t) - u_\phi(X_t,t)\big\|_2^2\Big].
\end{align}
After training, new samples are generated by simulating the ODE \eqref{eq:fm_ode} from $t=0$ to $t=T$, replacing the unknown velocity $u(x,t)$ with its learned estimator $u_\phi(x,t)$.


\textbf{Diffusion Policy and Flow Policy.}
A diffusion policy \cite{chi2025diffusion} is a policy induced by a diffusion model.
Let $s_\phi(x,t)$ be a parameterized score network, and let
$\{p_t^\phi\}_{t\in[0,T]}$ denote the marginal distributions obtained by
simulating the reverse-time SDE \eqref{eq:reverse_sde} with the score
$\nabla_x \log p_t(x)$ replaced by $s_\phi(x,t)$.
The diffusion policy $\pi_\phi$ is defined as the distribution over actions induced at the terminal time of the denoising process,
$\pi_\phi(\cdot|o) = p_0^\phi(\cdot|o)$.
A flow policy \cite{black2024pi,intelligence2025pi} is a policy induced by a flow matching model.
Let $u_\phi(x,t)$ be a parameterized velocity field, and let
$\{p_t^\phi\}_{t\in[0,T]}$ denote the marginals obtained by simulating the ODE \eqref{eq:fm_ode} with $u(x,t)$ replaced by $u_\phi(x,t)$.
The flow policy $\pi_\phi$ is defined as the action distribution at the terminal time of the flow, i.e.,
$\pi_\phi(\cdot|o)=p_T^\phi(\cdot|o)$.

For both diffusion and flow policies, evaluating the likelihood $\pi_\phi(a|o)$ is expensive.
As a result, the probability ratio
$\pi_\phi(a|o)/\pi_{\mathrm{old}}(a|o)$ in trust-region policy gradient methods (e.g., PPO, GRPO) is intractable.
This limits applying these methods to improve diffusion and flow policies.
DoobPO overcomes this obstacle and enables stable policy gradient style optimization for generative policies.